\documentclass[11pt]{article}

\usepackage[a4paper,margin=2.6cm]{geometry}
\usepackage{amsmath,amssymb}
\usepackage{hyperref}
\usepackage{xcolor}

\newcommand{\dd}{\mathrm{d}}
\newcommand{\ii}{\mathrm{i}}
\newcommand{\eps}{\epsilon}
\newcommand{\AMFlow}{\textsc{AMFlow}}
\newcommand{\WL}[1]{\texttt{#1}}

\title{AMFlow Project Log: Bubble, One-Loop, and Two-Loop Kernels}
\author{}
\date{\today}

\begin{document}
\maketitle

\section{Scope}

This note documents the current AMFlow sandbox project. The goal is to build
up confidence step by step before attempting the two-loop delta-constrained
counterterm \(I_C(x,y)\). The present project contains:
\begin{enumerate}
  \item a working ordinary one-loop bubble sanity check;
  \item a direct cut-linear one-loop collinear-kernel attempt, which currently
        fails inside AMFlow;
  \item a working uncut linear-propagator GaugeLink setup;
  \item a discontinuity reconstruction of the one-loop delta constraint from
        two uncut linear-propagator runs;
  \item a one-loop auxiliary-denominator test where the extra denominator is
        included in the family but requested with exponent zero;
  \item a first two-loop GaugeLink setup that reconstructs the two delta
        constraints from four uncut sign-prescription pieces.
\end{enumerate}
The two-loop setup is still at the first low-order AMFlow-test stage. The
purpose is to check whether AMFlow can reduce and evaluate the four uncut
two-loop GaugeLink pieces before increasing the requested Laurent order.

\section{Directories and External Setup}

The project directory is
\begin{verbatim}
/Users/FranzkeMM/Library/Mobile Documents/com~apple~CloudDocs/Documents/
  Dokumente privat/Studium/ETH Zürich/Master Thesis/amflow-project
\end{verbatim}
The AMFlow package is kept outside iCloud:
\begin{verbatim}
/Users/FranzkeMM/physics/tools/AMFlow
\end{verbatim}
The helper file
\begin{verbatim}
/Users/FranzkeMM/physics/tools/amflow-paths.wl
\end{verbatim}
adds the relevant paths for FiniteFlow, LiteIBP, LiteRed, and the FiniteFlow
dynamic library. A separate AMFlow-specific fix was also required: the reducer
file
\begin{verbatim}
AMFlow/ibp_interface/FiniteFlow+LiteRed/install.m
\end{verbatim}
must set
\begin{verbatim}
$FFPath
$FFLibraryPath
$LiteIBPPath
$LiteRedPath
\end{verbatim}
directly, because AMFlow's dependency check does not rely only on
\(\$\)\WL{Path}. Finally, \WL{fflowmlink.dylib} needed an rpath pointing to the
FiniteFlow build directory so that it can find \WL{libfflow.dylib}.

\section{Project Structure}

The active project files are:
\begin{itemize}
  \item \WL{config/LoadAMFlow.wl}: sets absolute project paths, loads
        \WL{amflow-paths.wl}, then loads AMFlow.
  \item \WL{config/AMFlowOptions.wl}: currently only sets the thread count.
  \item \WL{families/BubbleFamily.wl}: ordinary one-loop massive bubble.
  \item \WL{targets/RunBubble.wl}: runs the bubble sanity check.
  \item \WL{families/OneLoopKernelDirectCutFamily.wl}: direct cut-linear attempt.
  \item \WL{targets/RunOneLoopKernelDirectCut.wl}: runs the direct cut-linear attempt.
  \item \WL{families/OneLoopKernelUncutPlusFamily.wl}: uncut \(+L\) linear
        denominator family.
  \item \WL{targets/RunOneLoopKernelUncutPlus.wl}: runs the \(+L\) GaugeLink
        integral.
  \item \WL{families/OneLoopKernelUncutMinusFamily.wl}: uncut \(-L\) linear
        denominator family.
  \item \WL{targets/RunOneLoopKernelUncutMinus.wl}: runs the \(-L\) GaugeLink
        integral.
  \item \WL{targets/CompareOneLoop.wl}: computes \(F_+\),
        computes \(J_-\), forms the discontinuity, and compares to the
        analytic result in one fresh run.
  \item \WL{checks/CompareOneLoopFromFiles.wl}: reconstructs the
        delta constraint from previously exported uncut result files and
        compares coefficient by coefficient. This is convenient, but it can
        read stale files if the kinematics were changed.
  \item \WL{families/OneLoopKernelAuxDenFamilies.wl}: one-loop \(+L\) and
        \(-L\) families with one extra auxiliary denominator.
  \item \WL{targets/RunOneLoopKernelAuxDenPlus.wl}: runs the auxiliary
        denominator \(+L\) piece with the auxiliary exponent set to zero.
  \item \WL{targets/RunOneLoopKernelAuxDenMinus.wl}: runs the auxiliary
        denominator \(-L\) piece with the auxiliary exponent set to zero.
  \item \WL{targets/CompareOneLoopAuxDen.wl}: computes both auxiliary
        denominator pieces and compares to the analytic one-loop kernel.
  \item \WL{checks/CompareOneLoopAuxDenFromFiles.wl}: post-processes the
        exported auxiliary denominator pieces.
  \item \WL{checks/CompareOneLoopDirectCut.wl}: legacy comparison script for
        the direct cut-linear attempt.
  \item \WL{families/TwoLoopKernelUncutFamilies.wl}: four two-loop uncut
        GaugeLink families, labelled \WL{PP}, \WL{PM}, \WL{MP}, and \WL{MM}.
  \item \WL{targets/TwoLoopKernelTargetTools.wl}: shared helper functions for
        the two-loop targets.
  \item \WL{targets/RunTwoLoopKernelUncutPP.wl}: runs the \WL{PP} two-loop
        sign piece alone.
  \item \WL{targets/RunTwoLoopKernelUncutPM.wl}: runs the \WL{PM} two-loop
        sign piece alone.
  \item \WL{targets/RunTwoLoopKernelUncutMP.wl}: runs the \WL{MP} two-loop
        sign piece alone.
  \item \WL{targets/RunTwoLoopKernelUncutMM.wl}: runs the \WL{MM} two-loop
        sign piece alone.
  \item \WL{targets/CompareTwoLoop.wl}: computes all four two-loop pieces,
        forms the double discontinuity, and compares to the closed form.
  \item \WL{checks/CompareTwoLoopFromFiles.wl}: post-processes previously
        exported two-loop sign pieces.
\end{itemize}
The shell entry points are:
\begin{verbatim}
./run.sh bubble
./run.sh oneloop-kernel-direct-cut
./run.sh oneloop-kernel-uncut-plus
./run.sh oneloop-kernel-uncut-minus
./run.sh compare-oneloop
./run.sh compare-oneloop-from-files
./run.sh oneloop-kernel-aux-plus
./run.sh oneloop-kernel-aux-minus
./run.sh compare-oneloop-aux
./run.sh compare-oneloop-aux-from-files
./run.sh compare-oneloop-direct-cut
./run.sh twoloop-kernel-uncut-pp
./run.sh twoloop-kernel-uncut-pm
./run.sh twoloop-kernel-uncut-mp
./run.sh twoloop-kernel-uncut-mm
./run.sh compare-twoloop
./run.sh compare-twoloop-from-files
\end{verbatim}

\subsection*{Target Reference}

\begin{center}
\begin{tabular}{p{0.31\textwidth}p{0.23\textwidth}p{0.36\textwidth}}
\hline
\textbf{Target} & \textbf{Role} & \textbf{When to use it} \\
\hline
\WL{bubble}
& Ordinary one-loop bubble sanity check.
& Use this to verify that AMFlow, Mathematica subprocesses, the reducer, and
  the dynamic libraries still work. It is not a collinear-kernel test. \\
\WL{oneloop-kernel-direct-cut}
& Direct cut-linear attempt with \(L=n\cdot l+x\) marked as cut.
& Diagnostic only. This currently fails in AMFlow's cut-linear/boundary
  handling and is kept to document the failed direct route. \\
\WL{oneloop-kernel-uncut-plus}
& Computes \(F_+\), the uncut GaugeLink integral with denominator
  \(L+\ii0\).
& Diagnostic/cache target. Useful for inspecting the \(+L\) piece alone. By
  itself it is not the delta-constrained kernel. \\
\WL{oneloop-kernel-uncut-minus}
& Computes \(J_-\), the uncut GaugeLink integral with denominator
  \(-L+\ii0\).
& Diagnostic/cache target. Useful for inspecting the \(-L\) piece alone. By
  itself it is not the delta-constrained kernel. \\
\WL{compare-oneloop}
& Fresh full one-loop check. Computes \(F_+\), computes \(J_-\), forms
  \((-J_--F_+)/(2\pi\ii)\), and compares to the analytic kernel.
& Preferred target for the actual one-loop physics sanity check. It avoids
  stale result files because both AMFlow computations and the comparison happen
  in the same run. \\
\WL{compare-oneloop-from-files}
& Reads previously exported \(F_+\) and \(J_-\) result files, forms the
  discontinuity, and compares to the analytic expression.
& Quick post-processing only. Use with care: it can compare stale data if the
  result files were generated with old kinematics or old \(\eps\)-order. \\
\WL{oneloop-kernel-aux-plus}
& Computes \(F_+\) with one extra auxiliary denominator at exponent zero.
& Diagnostic target for testing whether adding a zero-power denominator changes
  the one-loop GaugeLink result. \\
\WL{oneloop-kernel-aux-minus}
& Computes \(J_-\) with one extra auxiliary denominator at exponent zero.
& Diagnostic target paired with \WL{oneloop-kernel-aux-plus}. \\
\WL{compare-oneloop-aux}
& Fresh auxiliary-denominator check. Computes both pieces, forms
  \((-J_--F_+)/(2\pi\ii)\), and compares to the analytic kernel.
& Preferred target for testing the zero-power denominator mechanism without
  relying on old result files. \\
\WL{compare-oneloop-aux-from-files}
& Reads previously exported auxiliary \(F_+\) and \(J_-\) result files and
  compares their discontinuity to the analytic expression.
& Quick post-processing only. Use only after rerunning both auxiliary pieces
  with the current setup. \\
\WL{compare-oneloop-direct-cut}
& Legacy comparison script for the direct cut-linear result.
& Not part of the recommended workflow at present, because the direct
  cut-linear AMFlow run does not currently produce the needed result. \\
\WL{twoloop-kernel-uncut-pp}
& Computes the two-loop \WL{PP} GaugeLink piece.
& First two-loop setup test. It checks one uncut sign piece without forming
  the full double discontinuity. \\
\WL{twoloop-kernel-uncut-pm}
& Computes the two-loop \WL{PM} GaugeLink piece.
& Diagnostic/cache target for the mixed \(L_x+\ii0\), \(-L_y+\ii0\) piece. \\
\WL{twoloop-kernel-uncut-mp}
& Computes the two-loop \WL{MP} GaugeLink piece.
& Diagnostic/cache target for the mixed \(-L_x+\ii0\), \(L_y+\ii0\) piece. \\
\WL{twoloop-kernel-uncut-mm}
& Computes the two-loop \WL{MM} GaugeLink piece.
& Diagnostic/cache target for the \(-L_x+\ii0\), \(-L_y+\ii0\) piece. \\
\WL{compare-twoloop}
& Fresh two-loop check. Computes all four sign pieces and compares the double
  discontinuity to the closed form.
& Preferred target once the individual pieces reduce successfully. It avoids
  stale exported files. \\
\WL{compare-twoloop-from-files}
& Reads exported two-loop sign pieces and forms the same comparison.
& Quick post-processing only. Use with care after changing kinematics or
  \(\eps\)-order. \\
\hline
\end{tabular}
\end{center}

\section{Stage 1: Bubble Sanity Check}

The bubble family is
\begin{equation}
  B =
  \int \frac{\dd^d l}{\ii\pi^{d/2}}\,
  \frac{1}{(l^2-m^2)((l+p)^2-m^2)}
\end{equation}
at the Euclidean point
\begin{equation}
  p^2=s=-1,\qquad m^2=\frac{1}{10}.
\end{equation}
This is not a physics target; it is only a sanity check that AMFlow, the
reducer, Mathematica subprocesses, and the dynamic libraries work.

The analytic calculation is also simple. Introduce a Feynman parameter,
\begin{equation}
  \frac{1}{D_1D_2}
  =
  \int_0^1 \dd z\,
  \frac{1}{\bigl[(l+zp)^2-m^2+z(1-z)p^2\bigr]^2}.
\end{equation}
After shifting the loop momentum and using \(d=4-2\eps\), the normalized
one-loop integral becomes
\begin{equation}
  B
  =
  \Gamma(\eps)
  \int_0^1\dd z\,
  \left[m^2-z(1-z)p^2\right]^{-\eps}.
\end{equation}
For the bubble sanity point \(p^2=-1\), \(m^2=1/10\),
\begin{equation}
  B
  =
  \Gamma(\eps)
  \int_0^1\dd z\,
  \left[\frac{1}{10}+z(1-z)\right]^{-\eps}.
\end{equation}
Expanding,
\begin{align}
  \Gamma(\eps)
  &=
  \frac{1}{\eps}-\gamma_E+\mathcal{O}(\eps),
  \\
  \left[\frac{1}{10}+z(1-z)\right]^{-\eps}
  &=
  1-\eps\log\!\left[\frac{1}{10}+z(1-z)\right]
  +\mathcal{O}(\eps^2),
\end{align}
so
\begin{equation}
  B
  =
  \frac{1}{\eps}
  -
  \gamma_E
  -
  \int_0^1\dd z\,
  \log\!\left[\frac{1}{10}+z(1-z)\right]
  +\mathcal{O}(\eps).
\end{equation}
The parameter integral can be evaluated as
\begin{equation}
  \int_0^1\dd z\,
  \log\!\left[a+z(1-z)\right]
  =
  \log a - 2
  +2\sqrt{1+4a}\,
  \operatorname{artanh}\!\left(\frac{1}{\sqrt{1+4a}}\right),
\end{equation}
with \(a=1/10\). Therefore
\begin{equation}
  B
  =
  \frac{1}{\eps}
  +0.7934919436865079827886470731\ldots
  +\mathcal{O}(\eps).
\end{equation}
This matches the AMFlow output one-to-one:
\begin{equation}
  B_{\AMFlow}
  =
  \frac{1}{\eps}
  +0.7934919436865079827886470731\ldots
  +\mathcal{O}(\eps).
\end{equation}

\section{Stage 2: Target One-Loop Kernel}

The one-loop differential collinear kernel to be tested is
\begin{equation}
  I_c(x,m,M)
  =
  \int \frac{\dd^d l}{\ii \pi^{d/2}}\,
  \frac{\delta(x+n\cdot l)}
       {(l^2-m^2)((l+p)^2-M^2)},
  \qquad d=4-2\eps ,
\end{equation}
with
\begin{equation}
  n^2=0,\qquad p\cdot n=1,\qquad p^2=0.
\end{equation}
The symbol \WL{eta} is avoided because AMFlow reserves it for the auxiliary
mass-flow variable. The light-like vector is called \WL{n}.

The AMFlow denominator convention is
\begin{equation}
  D_1=l^2-m^2,\qquad
  D_2=(l+p)^2-M^2,\qquad
  L=n\cdot l+x.
\end{equation}
The active numerical point is
\begin{equation}
  x=\frac13,\qquad m^2=\frac{1}{10},\qquad M^2=\frac15,\qquad p^2=0.
\end{equation}

\section{Stage 3: Direct Cut-Linear Attempt}

The most literal AMFlow encoding is to include \(L=n\cdot l+x\) as the third
propagator and set
\begin{equation}
  \WL{AMFlowInfo["Cut"] = \{0,0,1\}} .
\end{equation}
This is implemented in \WL{OneLoopKernelDirectCutFamily.wl}. It was useful as a
diagnostic, but it is not currently the working route:
\begin{itemize}
  \item ordinary \WL{SolveIntegrals} fails immediately because the standard
        path tries to rewrite all denominators as squared denominators;
  \item \WL{SolveIntegralsGaugeLink} accepts the linear denominator, but with
        the cut flag it fails later at an ending-system/boundary stage.
\end{itemize}
Thus the syntax \(n\,l+x\) is not the fundamental issue. The obstacle is direct
automatic treatment of a cut linear denominator.

\section{Stage 4: GaugeLink Linear-Propagator Tests}

AMFlow has a special function \WL{SolveIntegralsGaugeLink} for ordinary
linear propagators. It implements the strategy described in the AMFlow linear
propagator paper: introduce an auxiliary quadratic representation for a linear
denominator, solve differential equations in the auxiliary parameter, and take
the relevant limit.

We first tested the uncut integral
\begin{equation}
  F_+(x)
  =
  \int \frac{\dd^d l}{\ii \pi^{d/2}}\,
  \frac{1}{D_1D_2(L+\ii0)}.
\end{equation}
This has no direct delta-constraint meaning by itself, but it verifies that
AMFlow can handle the linear denominator machinery. The test succeeded.

The \(+L\) and \(-L\) families are kept as separate AMFlow family files because
AMFlow's reduction tables, cache directories, and integral heads are tied to a
single family definition. This separation was also useful during debugging.
However, a physics comparison should not rely on manually combining result
files from different runs, because those files may have been generated with
older kinematics. The preferred target is therefore
\begin{verbatim}
./run.sh compare-oneloop
\end{verbatim}
which performs both GaugeLink calculations and the comparison in one fresh
Mathematica session.

\section{Stage 5: Discontinuity Reconstruction}

Rather than cutting \(L\) directly in AMFlow, we reconstruct the delta function
from two ordinary linear-propagator integrals. The identity used is
\begin{equation}
  \delta(L)
  =
  \frac{1}{2\pi \ii}
  \left(
    \frac{1}{L-\ii0}
    -
    \frac{1}{L+\ii0}
  \right).
\end{equation}
The second AMFlow run uses the denominator \(-L\):
\begin{equation}
  J_-(x)
  =
  \int \frac{\dd^d l}{\ii \pi^{d/2}}\,
  \frac{1}{D_1D_2(-L+\ii0)}.
\end{equation}
Since
\begin{equation}
  \frac{1}{L-\ii0}=-\frac{1}{-L+\ii0},
\end{equation}
the physical discontinuity candidate is
\begin{equation}
  I_c^{\mathrm{disc}}(x,m,M)
  =
  \frac{-J_-(x)-F_+(x)}{2\pi\ii}.
\end{equation}
The comparison script also prints sign variants, but the expression above is
the intended one.

\section{Analytic Formula}

For the routing used here,
\begin{equation}
  I_c(x,m,M)
  =
  \frac{\Gamma(1+\eps)}{\eps}\,
  A^{-\eps}\,
  \Theta(0\le x\le 1),
\end{equation}
where
\begin{equation}
  A=(1-x)m^2+xM^2-x(1-x)p^2.
\end{equation}
With the intended light-like value \(p^2=0\),
\begin{equation}
  A=(1-x)m^2+xM^2.
\end{equation}
At the active point,
\begin{equation}
  A=\frac23\frac{1}{10}+\frac13\frac15=\frac{2}{15}.
\end{equation}
Therefore the expected expansion is
\begin{equation}
  \frac{\Gamma(1+\eps)}{\eps}
  \left(\frac{2}{15}\right)^{-\eps}.
\end{equation}
The finite coefficient is
\begin{equation}
  \log\!\left(\frac{15}{2}\right)-\gamma_E.
\end{equation}

\section{Important Correction: \(p^2\)}

An earlier one-loop run accidentally used \(s=p^2=-1\). This came from the
Euclidean bubble test and was not the intended collinear kinematics. With
\(s=-1\), the discontinuity check correctly matched the finite coefficient for
\begin{equation}
  A=(1-x)m^2+xM^2-x(1-x)s=\frac{16}{45},
\end{equation}
which was a useful sign/normalization check, but not the desired physics point.

The three one-loop family files now use
\begin{equation}
  s\to 0.
\end{equation}
Consequently, one-loop result files produced before this correction are stale.

\section{Current High-Order Test}

The preferred current high-order test is the single fresh-run target
\begin{verbatim}
./run.sh compare-oneloop
\end{verbatim}
It computes \(F_+\), computes \(J_-\), forms
\((-J_--F_+)/(2\pi\ii)\), and compares directly to the analytic expression.

The older split workflow is still available for diagnostics:
\begin{verbatim}
./run.sh oneloop-kernel-uncut-plus
./run.sh oneloop-kernel-uncut-minus
./run.sh compare-oneloop-from-files
\end{verbatim}
but the last command reads exported files and can therefore compare stale data
if one of the two calculations was not rerun after a configuration change.

The two GaugeLink target scripts have been set to
\begin{equation}
  \WL{epsOrder = 5}.
\end{equation}
For \WL{SolveIntegralsGaugeLink}, this is necessary because AMFlow internally
uses a one-loop leading power \(-2\), so the returned expansion reaches
\(\eps^{\WL{epsOrder}-2}\). Thus \WL{epsOrder = 5} returns coefficients
through \(\eps^3\). The comparison script expands the analytic result through
\(\eps^3\) and compares powers
\begin{equation}
  \eps^{-1},\eps^0,\eps^1,\eps^2,\eps^3.
\end{equation}
The comparison is coefficient-wise:
\begin{equation}
  c_p^{\mathrm{AMFlow}}
  =
  [\eps^p]\,
  \frac{-J_--F_+}{2\pi\ii},
  \qquad
  c_p^{\mathrm{ana}}
  =
  [\eps^p]\,
  \frac{\Gamma(1+\eps)}{\eps}A^{-\eps}.
\end{equation}
For each \(p\), the script prints
\begin{equation}
  c_p^{\mathrm{AMFlow}},\qquad
  c_p^{\mathrm{ana}},\qquad
  c_p^{\mathrm{AMFlow}}-c_p^{\mathrm{ana}},\qquad
  \frac{c_p^{\mathrm{AMFlow}}}{c_p^{\mathrm{ana}}}.
\end{equation}

\section{Stage 5b: One-Loop Auxiliary-Denominator Test}

The two-loop family needs an extra denominator to complete the scalar-product
space. It is then requested with exponent zero. Before trusting this mechanism
in the expensive two-loop run, the same idea is tested in the one-loop setup.
The auxiliary one-loop families use
\begin{equation}
  D_1=l^2-m^2,\qquad
  D_2=(l+p)^2-M^2,\qquad
  D_3=\pm(n\cdot l+x),\qquad
  D_4=(l-p)^2.
\end{equation}
The AMFlow targets are
\begin{equation}
  \WL{j[oneloopkernelauxplus, 1, 1, 1, 0]},
  \qquad
  \WL{j[oneloopkernelauxminus, 1, 1, 1, 0]}.
\end{equation}
Thus \(D_4\) is present in the family definition but absent from the integral.
The physical one-loop integral should therefore be unchanged.

The fresh target is
\begin{verbatim}
./run.sh compare-oneloop-aux
\end{verbatim}
It computes the auxiliary \(F_+\) and \(J_-\) pieces, forms the same
discontinuity,
\begin{equation}
  I_{c,\mathrm{aux}}^{\mathrm{disc}}
  =
  \frac{-J_-^{\mathrm{aux}}-F_+^{\mathrm{aux}}}{2\pi\ii},
\end{equation}
and compares it coefficient by coefficient to
\begin{equation}
  \frac{\Gamma(1+\eps)}{\eps}A^{-\eps},
  \qquad
  A=(1-x)m^2+xM^2-x(1-x)p^2.
\end{equation}
The split workflow is also available:
\begin{verbatim}
./run.sh oneloop-kernel-aux-plus
./run.sh oneloop-kernel-aux-minus
./run.sh compare-oneloop-aux-from-files
\end{verbatim}
The file-based command is only post-processing. It must not be trusted after a
change of kinematics or \(\eps\)-order unless both auxiliary pieces were rerun.

\section{Stage 6: Two-Loop GaugeLink Setup}

The two-loop target is the double-collinear kernel from
\WL{derivations/snippets/dckernel/main.tex}:
\begin{equation}
  \mathcal I(x,y)=
  \int \frac{\dd^d k}{\ii\pi^{d/2}}
       \frac{\dd^d l}{\ii\pi^{d/2}}\,
  \frac{\delta(x-l\cdot n)\delta(y-k\cdot n)}
       {(l^2-m_l^2)((l-p)^2-M_l^2)(k^2-m_k^2)((l+k-p)^2-M_k^2)} .
\end{equation}
The AMFlow setup uses
\begin{equation}
  n^2=0,\qquad p\cdot n=1,
\end{equation}
and the two linear denominators
\begin{equation}
  L_x=x-n\cdot l,\qquad L_y=y-n\cdot k.
\end{equation}
As in the one-loop test, the direct cut-linear route is not used. Instead,
four uncut GaugeLink pieces are introduced:
\begin{align}
  I_{PP} &: \frac{1}{(L_x+\ii0)(L_y+\ii0)},\\
  I_{PM} &: \frac{1}{(L_x+\ii0)(-L_y+\ii0)},\\
  I_{MP} &: \frac{1}{(-L_x+\ii0)(L_y+\ii0)},\\
  I_{MM} &: \frac{1}{(-L_x+\ii0)(-L_y+\ii0)}.
\end{align}
The double delta constraint is reconstructed from
\begin{equation}
  \delta(L_x)\delta(L_y)
  =
  \frac{1}{(2\pi\ii)^2}
  \left(
    \frac{1}{L_x-\ii0}-\frac{1}{L_x+\ii0}
  \right)
  \left(
    \frac{1}{L_y-\ii0}-\frac{1}{L_y+\ii0}
  \right).
\end{equation}
Using
\begin{equation}
  \frac{1}{L-\ii0}=-\frac{1}{-L+\ii0},
\end{equation}
the implemented double-discontinuity candidate is
\begin{equation}
  \mathcal I_{\mathrm{disc}}
  =
  \frac{I_{PP}+I_{PM}+I_{MP}+I_{MM}}{(2\pi\ii)^2}.
\end{equation}

The first numerical point is now a light-like point chosen to avoid branch
phases while debugging the discontinuity normalization:
\begin{equation}
  x=\frac{3}{10},\qquad y=\frac14,\qquad p^2=0,
\end{equation}
and
\begin{equation}
  m_l^2=\frac{49}{100},\qquad
  M_l^2=1,\qquad
  m_k^2=\frac14,\qquad
  M_k^2=\frac{81}{100}.
\end{equation}
This differs from the older Euclidean Python-check point. Both \(\Delta_0\)
and \(\Delta_1\) are positive at this point, so the first AMFlow comparison is
not mixed with a branch-cut phase test.
The analytic comparison uses
\begin{equation}
  \mathcal I_{\mathrm{closed}}
  =
  \frac{\Gamma(\eps)^2}{1-x}
  [\Delta_0-\ii0]^{-\eps}
  [\Delta_1-\ii0]^{-\eps},
\end{equation}
with
\begin{equation}
  \Delta_0=(1-x)m_l^2+xM_l^2-x(1-x)p^2,
\end{equation}
and
\begin{equation}
  \Delta_1=
  \left(1-\lambda\right)m_k^2+\lambda M_k^2
  -\lambda(1-\lambda)M_l^2
  -\lambda(1-\lambda x)p^2,
  \qquad
  \lambda=\frac{y}{1-x}.
\end{equation}
At the current stage the two-loop targets use
\begin{equation}
  \WL{precisionGoal = 10},\qquad \WL{epsOrder = 4},
\end{equation}
and compare the powers
\begin{equation}
  \eps^{-2},\eps^{-1},\eps^0.
\end{equation}
This is still a modest run, but it should include the finite coefficient.

For quick smoke tests, the same scripts can be run at lower precision while
keeping the same \(\eps\)-order. The predefined shortcut
\begin{verbatim}
./run.sh compare-twoloop-quick
\end{verbatim}
sets
\begin{equation}
  \WL{precisionGoal = 6},\qquad \WL{epsOrder = 4},
\end{equation}
unless these values are overridden by environment variables. This should be
used only to catch setup mistakes and coefficient-level problems at rough
precision. A slightly more accurate check is, for example,
\begin{verbatim}
AMFLOW_PRECISION_GOAL=8 ./run.sh compare-twoloop-quick
\end{verbatim}
The single-piece shortcut
\begin{verbatim}
./run.sh twoloop-kernel-uncut-pp-quick
\end{verbatim}
is useful when testing whether one representative family still reduces.

The first completed run with the older \(p^2=-3\) point produced a leading
coefficient exactly one half of the closed-form value. The comparison scripts
therefore now print an inferred leading-pole normalization,
\begin{equation}
  \frac{[\eps^{-2}]\mathcal I_{\mathrm{disc}}}
       {[\eps^{-2}]\mathcal I_{\mathrm{closed}}},
\end{equation}
and expose a variable \WL{doubleCutNormalization}. It is initially set to
\(1\). If the light-like rerun again gives a constant factor \(1/2\), set
\WL{doubleCutNormalization = 1/2} in the comparison script and rerun only the
post-processing comparison. No AMFlow recomputation is needed for that
normalization check.

The recommended debugging order is:
\begin{verbatim}
./run.sh twoloop-kernel-uncut-pp
\end{verbatim}
If that succeeds, try the remaining pieces or the fresh full check:
\begin{verbatim}
./run.sh compare-twoloop
\end{verbatim}
For quick post-processing after all four sign pieces have already been
exported, use:
\begin{verbatim}
./run.sh compare-twoloop-from-files
\end{verbatim}
The file-based comparison can read stale data, so it should not be used after
changing the kinematic point or \(\eps\)-order unless all four pieces were
rerun.


\end{document}
